\section{Mô hình hệ thống và kiến trúc}
\label{sec:system-model}

\subsection{Mô hình đe dọa và phạm vi}
\label{subsec:threat-model}

Giải pháp là một công cụ hậu kiểm nhật ký cho ứng dụng kiểu ngân hàng. Nó tách
hai trách nhiệm: mô hình RBAC biểu diễn quan hệ phân quyền, còn pipeline phân
tích dùng trạng thái RBAC và log đã ghi để phát hiện rủi ro. Hệ thống giả định
một dịch vụ nguồn đã ghi sự kiện theo sơ đồ quy định và nhà vận hành cung cấp cơ
sở dữ liệu RBAC cùng cấu hình đã được kiểm soát. Việc thu thập, chuẩn hóa và
phân tích nhật ký là một phần quan trọng của điều tra sự cố
\citep{Kent2006SP80092}; ở đây, các bước đó được thu gọn thành một pipeline
nhỏ có thể tái lập, không phải kiến trúc quản lý log đầy đủ.

Kẻ tấn công trong phạm vi có thể tạo ra các hành vi để lại dấu vết trong nhật
ký: gửi request chứa chỉ báo SQL injection hoặc lặp lại xác thực thất bại để dò
đoán mật khẩu. Kiểm tra event truy cập/phân quyền không phù hợp với RBAC là
nguồn bằng chứng hỗ trợ. Các kịch bản này tương ứng với dấu hiệu khai thác ứng
dụng công khai và brute force ở mức quan sát được
\citep{OWASPSQLi,MITRET1190,MITRET1110}. Mục tiêu là tạo cảnh báo có evidence
và ưu tiên chúng để hỗ trợ triage, không phải xác nhận một cuộc tấn công thành
công.

\subsection{Đầu vào, đầu ra và ranh giới tin cậy}
\label{subsec:trust-boundary}

Đầu vào không tin cậy là tệp nhật ký CSV hoặc JSON; các trường như user, IP,
request và status được parser kiểm tra trước khi chuyển thành event. Nhãn
\texttt{expected\_label} trong các tệp demo chỉ là nhãn đánh giá và không đi
vào quyết định phát hiện. CSDL SQLite chứa trạng thái người dùng--vai trò--quyền
và tệp TOML chứa ngưỡng, cửa sổ thời gian và điểm cộng; hai thành phần này thuộc
ranh giới tin cậy của lần chạy. Ranh giới đó cũng bao gồm mã nguồn dịch vụ và thư
mục xuất report. Do đó, kết quả chỉ đáng tin cậy khi RBAC seed, cấu hình và
nguồn log được nhà vận hành bảo vệ khỏi sửa đổi.

Đầu ra gồm alert theo luật hoặc theo ngữ cảnh, metadata của lần chạy và, khi
bật context-risk, context finding cùng incident. Mỗi alert giữ định danh event
và evidence; mỗi incident giữ các alert/event IDs, evidence, điểm và hai mốc
thời gian của đoạn nhóm. Thiết kế này làm rõ đường truy vết từ report về input,
nhưng không thay thế cơ chế toàn vẹn, ký số hoặc lưu trữ bất biến của log.

\subsection{Phạm vi phát hiện và các điều không xử lý}
\label{subsec:scope-limitations}

Nguyên mẫu tập trung vào ba nhóm rủi ro trong nhật ký ứng dụng kiểu ngân hàng:

\begin{enumerate}
  \item Dấu hiệu SQL injection trong chuỗi yêu cầu.
  \item Đoán mật khẩu, được biểu diễn bằng nhiều lần đăng nhập thất bại.
  \item Truy cập trái phép, được biểu diễn bằng hành động bị RBAC hoặc trạng
  thái nhật ký từ chối.
\end{enumerate}

Hệ thống không xử lý dữ liệu ngân hàng thật, không thực thi quyết định cấp quyền
trực tuyến, không bảo vệ kênh thu thập log, không chống giả mạo/sửa đổi log và
không triển khai tương quan đa nguồn, điều tra pháp chứng hay phản ứng tự động.
Nó cũng không học mô hình hành vi hoặc đánh giá độ chính xác trên dữ liệu vận
hành. Đây là minh chứng ý tưởng theo hướng giảng dạy, nhằm làm cho mỗi cảnh báo
có thể truy vết đến event, luật và evidence; một triển khai vận hành sẽ cần các
kiểm soát rộng hơn \citep{Rose2020SP800207,JTF2020SP80053}.

\subsection{Kiến trúc}
\label{subsec:architecture}

Hình~\ref{fig:architecture} tóm tắt pipeline xử lý.

\begin{figure}[t]
  \centering
  \resizebox{\linewidth}{!}{% Luồng thực thi trong src/rbac_guard/service.py.
\begin{tikzpicture}[
  font=\footnotesize,
  >=Latex,
  node distance=7mm and 9mm,
  component/.style={draw, rounded corners=2pt, align=center, minimum height=12mm,
    text width=25mm, fill=blue!7},
  source/.style={component, fill=gray!12},
  store/.style={component, fill=orange!14},
  output/.style={component, fill=green!11},
  flow/.style={->, line width=.55pt},
  auxiliary/.style={->, dashed, line width=.5pt}
]
  \node[source] (logs) {Nhật ký CSV/JSON\\và cấu hình};
  \node[component, right=of logs] (parser) {Parser\\sự kiện hợp lệ};
  \node[store, above=of parser] (rbac) {CSDL SQLite\\RBAC};
  \node[component, right=of parser] (rules) {Bộ máy luật\\SQLi, mật khẩu, RBAC};
  \node[component, below=of rules] (context) {Profiler ngữ cảnh\\giờ, IP, tài nguyên, từ chối};
  \node[component, right=of rules] (scorer) {Bộ chấm điểm\\và mức độ};
  \node[output, right=of scorer] (reporter) {Reporter\\alerts, findings, incidents\\CSV/JSON};
  \node[source, below=of parser] (adapter) {CLI / Streamlit\\adapter};

  \draw[flow] (logs) -- (parser);
  \draw[flow] (parser) -- (rules);
  \draw[flow] (parser) -| (context);
  \draw[auxiliary] (rbac) -- (rules);
  \draw[flow] (rules) -- (scorer);
  \draw[flow] (context) -| (scorer);
  \draw[flow] (scorer) -- (reporter);
  \draw[auxiliary] (adapter) -- (parser);
\end{tikzpicture}
}
  \caption{Kiến trúc giải pháp. Đường chính gồm parser, bộ máy luật có đối
  chiếu RBAC, chấm điểm và reporter. Profiler ngữ cảnh cùng incident chỉ được
  kích hoạt trong chế độ context-risk tùy chọn.}
  \label{fig:architecture}
\end{figure}

Giao diện dòng lệnh và giao diện Streamlit là các adapter phía trên một dịch vụ
ứng dụng dùng chung. Đường xử lý chính phân tích nhật ký đầu vào, khởi tạo hai
luật phát hiện chính và luật đối chiếu RBAC, chấm điểm finding, rồi ghi alert và
metadata. Khi bật chế độ tùy chọn, dịch vụ mới tính tín hiệu ngữ cảnh và gom
alert thành incident.

\subsection{Mô hình dữ liệu}
\label{subsec:data-model}

Mỗi sự kiện gồm định danh, dấu thời gian, loại sự kiện, người dùng, IP nguồn,
tài nguyên, hành động, trạng thái, yêu cầu và chi tiết. Các dataset đánh giá
trong kho mã nguồn cung cấp thêm nhãn kỳ vọng; parser hiện yêu cầu trường này,
nhưng logic phát hiện không dùng nhãn để quyết định cảnh báo. Dữ liệu RBAC được
lưu trong các bảng SQLite cho người
dùng, vai trò, quyền, gán người dùng--vai trò và gán vai trò--quyền.

Mô hình đầu ra cốt lõi là alert: mỗi alert biểu diễn một finding theo luật, giữ
event IDs, rule ID, evidence, điểm rủi ro và severity. Khi bật context-risk,
hệ thống sinh thêm finding ngữ cảnh và incident; incident gom alert liên quan
theo người dùng, IP và cửa sổ thời gian, đồng thời chứa phần tóm tắt, điểm,
mức độ nghiêm trọng và chuỗi bằng chứng.
